% Statements about cited works are limited to what their abstracts state.
\section{Related work}
\label{sec:related}

Our statements about cited works rest on their abstracts and bibliographic records, not on a full reading; the Handbook of Computational Social Choice and the evaluation harness are cited as general references.

\paragraph{Fragility of multi-task benchmarks.}
\citet{dehghani2021benchmark} show, on several common benchmark setups, that the relative performance of methods
can change substantially with the choice of benchmark tasks (the ``benchmark lottery'').
Their perturbation is the choice of tasks by the evaluator; ours is a targeted change of one model's scores,
and our question is quantitative: how much change is needed, under which rule.
\citet{lanctot2023evaluating} view each task as a voter and bring social-choice rules to cross-task evaluation of agents and language models;
we use their viewpoint to select rank-based rules (Borda, Condorcet winner, plurality) next to score-based rules
(mean, median, minimum, geometric, harmonic, power and trimmed means), and we compare how costly it is to win under each.
Standard references on these voting rules are collected in the Handbook of Computational Social Choice \citep{brandt2016handbook}.

\paragraph{Strategic behaviour and contamination.}
\citet{singh2025leaderboard} document practices on the Chatbot Arena leaderboard (private testing of many variants, selective score disclosure, unequal sampling)
that distort rankings. Benchmark leakage and contamination, i.e.\ data related to an evaluation set being used for training,
is studied empirically by \citet{zhou2023dont}, who report that it can substantially boost scores,
and \citet{yang2023rethinking}, who show that rephrased test samples evade common decontamination and let a 13B model overfit a benchmark.
These works establish that targeted improvement of a benchmark's score is feasible and consequential;
they do not ask what the \emph{minimum} improvement for first place is, nor how it depends on how sub-dataset scores are aggregated.
Our cap scenarios (Section~\ref{sec:method}) are a modelling stand-in for what such data can achieve, not a measurement of it.
The closest listed work, \citet{chen2026leaderboard}, models benchmarking as a Stackelberg game between a benchmark designer and several developers
who can inflate observed scores by allocating resources to benchmark-specific improvements, and gives equilibrium results.
We differ in setting and method: a single target against fixed rivals, exact minimum-cost solutions per aggregation rule,
and empirical evaluation on public leaderboard data.

\paragraph{Minimum-cost manipulation in voting.}
Computational social choice studies the cost of changing an election's input so that a chosen candidate wins,
e.g.\ bribery, where an external actor pays voters to change their preferences \citep{faliszewski2009bribery}.
Minimum investment to rank first is the analogue for score matrices: the ``voters'' are tasks, the ``bribe'' is a raise of the target's score on a task,
and the cost is measured in score points or number of tasks touched. The setting differs in that the manipulated quantity is a cardinal score
bounded by a headroom cap, and the instances are real leaderboards with $n$ of order $6$--$57$ tasks.

\paragraph{Evaluation data.}
We use the HELM MMLU release \citep{liang2022holistic} over the 57 MMLU subjects \citep{hendrycks2021measuring},
and the Open LLM Leaderboard (tasks IFEval, BBH \citep{suzgun2023challenging}, MATH Lvl~5, GPQA, MuSR, MMLU-Pro), whose tasks are run with the Language Model Evaluation Harness \citep{gao2023harness}.
